% Evans 11.2, printed pp.582--595; English PDF pp.599--612.
% Chinese start: printed p.412 / scan p.430.
\section{Riemann's Problem}
\label{sec:11-2}
The initial data are $u_l$ for $x<0$ and $u_r$ for $x>0$.
All constructions below take place in a sufficiently small strictly
hyperbolic neighborhood of a fixed state. Distinguish the state-space
curves $R_k,S_k$ from their space-time realizations: the former
parameterize possible states, whereas the latter have slopes $x/t$.
\PDEAnnotation{标量化的适用范围}{简单波代入使用链式法则。
连续稀疏波可逐片验证再拼接；若标量 $w$ 跳跃，
$V(w)$ 未必满足系统 Rankine--Hugoniot 条件，不能直接搬用标量激波。}

\subsection{Simple waves and genuine nonlinearity}
Seek $u=V(w(x,t))$ with $V'=r_k(V)$. Substitution gives
\[
 u_t+DF(u)u_x=r_k(V)\{w_t+\lambda_k(V(w))w_x\}.
\]
Thus a smooth scalar solution of $w_t+f_k(w)_x=0$, where
$f_k'(s)=\lambda_k(V(s))$, produces a smooth system solution.
The integral curve $V$ through $z$ is the $k$-rarefaction curve $R_k(z)$.
\begin{definition}[Source \S11.2.1]
The $k$-field is genuinely nonlinear if
$\gamma_k=D\lambda_k\cdot r_k\ne0$ throughout the neighborhood;
it is linearly degenerate if $\gamma_k\equiv0$ there.
\end{definition}
For a genuinely nonlinear field orient the unit vector so that
$\gamma_k>0$. Then $f_k''=\gamma_k(V)>0$ locally. Reversing the
orientation changes the parameter but not the geometric wave.

\begin{theorem}[Centered rarefaction; source \S11.2.2, Theorem 1]
Suppose the $k$-field is genuinely nonlinear and $u_r=V(s_r)$,
$u_l=V(s_l)$ with $s_l<s_r$. On this short curve let
$L(s)=\lambda_k(V(s))$. Then the continuous integral solution is
\begin{equation}\label{eq:11-fan}
 u(x,t)=\begin{cases}
 u_l,&x/t\le L(s_l),\\
 V(L^{-1}(x/t)),&L(s_l)<x/t<L(s_r),\\
 u_r,&x/t\ge L(s_r).
 \end{cases}
\end{equation}
\end{theorem}
\begin{expandedproof}[PDE, interfaces, and the initial trace]
Since $L'=\gamma_k>0$, its inverse exists and is smooth. In the fan,
$U(\xi)=V(L^{-1}(\xi))$ satisfies
$(DF(U)-\xi I)U'=0$, so $u_t+DF(u)u_x=0$ for $t>0$.
On a strip $t\ge\delta>0$ the profile is piecewise $C^1$ and continuous
at both fan edges; therefore there is no interface measure in its
first derivatives. The fan and all differing initial states occupy
an interval of length at most $Ct$ near the origin. Boundedness gives
$\|u(\cdot,t)-g\|_{L^1(K)}\le C_Kt$ for each compact $K$.
Integration by parts above $\delta$ followed by $\delta\downarrow0$
proves the initial weak identity.
\end{expandedproof}
\PDEBodyDiagram{11.1}{Wave patterns in physical space}
{The left panel is a centered rarefaction: characteristic speeds increase
from $\lambda_l$ to $\lambda_r$. The middle panel shows a compressive
shock with $\lambda_r<\sigma<\lambda_l$; forward characteristics terminate
on it. The right panel is a contact, with equal characteristic and interface
speed. Only slopes and incidence are represented; these are schematic
panels for a single characteristic family.}{ch11-wave-patterns}

\subsection{The Hugoniot locus and its second-order contact}
The shock set at $z$ consists of states $v$ for which
$F(v)-F(z)=\sigma(v-z)$ for some common speed $\sigma$.
\begin{theorem}[Local shock curves; source \S11.2.3, Theorem 2]
For $F\in C^3$ strictly hyperbolic near $z$, the shock set is locally
the union of $m$ curves $S_k(z)$. In a signed unit-speed parameter
$a$, choose $S_k(0;z)=z$, $\partial_aS_k(0;z)=r_k(z)$.
Writing $v=S_k(a;z)$, the shock speed extends smoothly along this
branch to $\sigma_k(0;z)=\lambda_k(z)$, and
\begin{equation}\label{eq:11-shock-expansion}
 \partial_a\sigma_k(0;z)=\tfrac12\gamma_k(z),\qquad
 \sigma_k(a;z)=\frac{\lambda_k(v)+\lambda_k(z)}2+O(a^2).
\end{equation}
The shock and rarefaction curves have equal second derivatives at zero
when both use oriented arc length.
\end{theorem}
\PDEAnnotation{激波速度不是全域函数}{每一条 $S_k$ 上的速度可光滑
延拓到基点，但不同分支极限是不同的 $\lambda_k(z)$。
不能在所有分支的交点假定一个共同的光滑 $\sigma(v,z)$。}
\begin{expandedproof}[construction and all Taylor coefficients]
Define $\overline A(v,z)=\int_0^1DF(z+\theta(v-z))\,d\theta$.
The jump equation becomes
$(\overline A(v,z)-\sigma I)(v-z)=0$.
Near $v=z$ this averaged matrix is strictly hyperbolic; its real
simple eigenvalues and dual eigenvectors are smooth. Reality follows
by continuation of the real simple roots of the characteristic
polynomial, and the spectral gaps persist after shrinking the ball.
Let these left vectors be $\overline l_j(v,z)$. The $m-1$ equations
\[
 \overline l_j(v,z)\cdot(v-z)=0,\qquad j\ne k,
\]
 have derivative at $v=z$ with independent rows $l_j(z)$.
The implicit function theorem gives a curve tangent to $r_k(z)$.
Every nearby nontrivial jump vector is an eigenvector of the averaged
matrix, hence belongs to exactly one of these branches. The speed
on the branch is its $k$th eigenvalue; at the common base point there
are $m$ distinct branch limits, not a single joint speed function.

Set $r=r_k(z)$, $A=DF(z)$, $\lambda=\lambda_k(z)$ and
$H=D^2F(z)[r,r]$. Twice differentiating
$F(S(a))-F(z)=\sigma(a)(S(a)-z)$ at zero gives
\[
 (A-\lambda I)S''(0)=2\sigma'(0)r-H.
\]
Differentiate $DF\,r_k=\lambda_k r_k$ in direction $r$:
\[
 (A-\lambda I)Dr_k(z)r=\gamma_k(z)r-H.
\]
Left multiplication by $l_k$ yields $2\sigma'(0)=\gamma_k$.
Subtracting the equations shows that $S''(0)-Dr_k r$ is parallel
to $r$. Both terms are orthogonal to $r$, because both curves have
unit speed; hence they are equal. Taylor's formula now gives
$\lambda_k(S(a))=\lambda+\gamma_k a+O(a^2)$ and the stated speed
estimate. Local equivalence $|S(a)-z|\asymp|a|$ also permits the
remainder to be written $O(|v-z|^2)$.
\end{expandedproof}

\begin{theorem}[Contacts; source \S11.2.3, Theorem 3]
If the $k$-field is linearly degenerate on a neighborhood, its local
shock and rarefaction curves coincide and their speed is
$\lambda_k(z)=\lambda_k(v)$. Two such states are joined by a contact
at that speed.
\end{theorem}
\begin{expandedproof}[inclusion and local equality]
Along $V'=r_k(V)$, $\frac d{ds}\lambda_k(V)=0$.
Integration gives
$F(V(s))-F(z)=\int_0^s\lambda_k(z)V'(a)\,da
=\lambda_k(z)(V(s)-z)$.
Thus the rarefaction curve lies in the shock set. Its tangent places
it on branch $k$; the implicit-function construction gives the unique
local curve on that branch. The two germs coincide, and
\eqref{eq:11-rh} gives the contact solution. This is a local equality,
not a statement about arbitrary distant branches of the shock set.
\end{expandedproof}

\subsection{Admissible shocks and the local solver}
For a genuinely nonlinear field with $\gamma_k>0$, \eqref{eq:11-shock-expansion}
gives for small $a<0$
\[
 \lambda_k(S_k(a;z))<\sigma_k(a;z)<\lambda_k(z).
\]
These are the compressive Lax shocks. For $a>0$ the inequalities
reverse and give expansion shocks, which are rejected. Spectral
separation also gives
$\lambda_{k-1}(z)<\sigma_k<\lambda_{k+1}(S_k(a;z))$
when the neighboring indices exist. The strict inequalities follow
because the leading gaps have size $|\gamma_k a|/2$, whereas the
error has size $O(a^2)$.
\PDEMarginDiagram{11.M}{平均速度与压缩}
{本图取 Burgers 通量 $f(u)=u^2/2$、$u_l=1,u_r=0$。
激波速度恰为 $1/2$，两侧速度为 $1$ 和 $0$；
一般系统的平均公式只在小振幅下成立到二阶余项。}{ch11-speed-margin}

Define $T_k(a;z)$ as the rarefaction curve for $a\ge0$ and the
shock curve for $a\le0$, each with compatible oriented parameters.
For a linearly degenerate field use its contact curve for both signs.
The common value and tangent, as well as smooth base-state dependence,
make $T_k$ jointly $C^1$ near $(0,z)$; the second-order contact proved
above gives additional matching in the wave parameter.
\begin{theorem}[Local Riemann solver; source \S11.2.4, Theorem 4]
If each field is genuinely nonlinear or linearly degenerate on a
neighborhood of $u_l$, every sufficiently close $u_r$ can be joined
to $u_l$ by at most $m$ elementary waves, ordered by family.
The resulting integral solution is self-similar. Its small wave-strength
vector is unique within this construction.
\end{theorem}
\begin{expandedproof}[state matching, speed ordering, weak gluing]
Let $z_0=u_l$ and $z_k=T_k(a_k;z_{k-1})$. The endpoint map
$\mathcal E(a)=z_m$ is $C^1$, $\mathcal E(0)=u_l$, and
\[
 D\mathcal E(0)=(r_1(u_l),\ldots,r_m(u_l)).
\]
This is invertible. The inverse function theorem gives a unique small
$a$ with $\mathcal E(a)=u_r$. Choose disjoint intervals around the
ordered numbers $\lambda_k(u_l)$. For states and strengths sufficiently
small, every fan edge, shock speed, and contact speed in family $k$
remains in its own interval. Consequently the elementary waves do not
overlap in $\xi=x/t$, and constant intermediate states fill the gaps.
Integrate on each smooth sector: continuous fan boundaries have no
jump measure, and each shock or contact satisfies
\eqref{eq:11-rh}. All interior interface terms cancel. The difference
from the initial step is bounded and confined to $|x|\le Ct$, giving
the initial trace as in the rarefaction proof. The inverse theorem
proves uniqueness of the small parameters; it does not prove uniqueness
among all bounded weak or entropy solutions of the system.
\end{expandedproof}
\PDEAnnotation{波曲线的方向}{本稿始终写右状态 $z_k\in T_k(z_{k-1})$。
原书 §11.2.4 的式 (43) 与随后 (45) 的左右方向写法不一致；
这里以构造式 (45) 为准，明确每次从左状态走向右状态。}
\PDEBodyDiagram{11.2}{Matching states before placing waves}
{The endpoint map composes the family curves in state space. Its derivative
at zero consists of the independent eigenvectors. Disjoint speed intervals
then place those waves in increasing order in $x/t$. The unique small
strength vector is a local finite-dimensional conclusion; no general
system uniqueness theorem is inferred.}{ch11-local-solver}
